% =============================================================================
%  Documentação Técnica — Natalia Time v1.1.0
%  Compilar com: pdflatex documentacao_tecnica.tex  (2 passagens)
% =============================================================================
\documentclass[12pt,a4paper]{report}

% ── Codificação e idioma ──────────────────────────────────────────────────────
\usepackage[utf8]{inputenc}
\usepackage[T1]{fontenc}
\usepackage[brazil]{babel}

% ── Layout ───────────────────────────────────────────────────────────────────
\usepackage[top=2.5cm, bottom=2.5cm, left=2.8cm, right=2.8cm]{geometry}
\usepackage{fancyhdr}
\usepackage{titlesec}

% ── Matemática ───────────────────────────────────────────────────────────────
\usepackage{amsmath}
\usepackage{amssymb}

% ── Gráficos e cores ─────────────────────────────────────────────────────────
\usepackage{graphicx}
\usepackage[table]{xcolor}

% ── Tabelas ──────────────────────────────────────────────────────────────────
\usepackage{booktabs}
\usepackage{longtable}
\usepackage{array}

% ── Listas personalizadas ────────────────────────────────────────────────────
\usepackage{enumitem}

% ── Código-fonte ─────────────────────────────────────────────────────────────
\usepackage{listings}

% ── Hyperlinks ───────────────────────────────────────────────────────────────
\usepackage[hidelinks, colorlinks=true,
            linkcolor=blue!60!black,
            urlcolor=blue!60!black,
            citecolor=blue!60!black]{hyperref}

% ── Definição de cores ───────────────────────────────────────────────────────
\definecolor{codebg}{RGB}{245,245,245}
\definecolor{codeframe}{RGB}{200,200,200}
\definecolor{keyword}{RGB}{0,0,180}
\definecolor{comment}{RGB}{0,128,0}
\definecolor{string}{RGB}{163,21,21}
\definecolor{azul}{RGB}{30,100,200}
\definecolor{cinzatitulo}{RGB}{50,50,80}

% ── Estilo de código Python ───────────────────────────────────────────────────
\lstdefinestyle{python}{
  language=Python,
  backgroundcolor=\color{codebg},
  frame=single,
  framerule=0.6pt,
  rulecolor=\color{codeframe},
  basicstyle=\ttfamily\small,
  keywordstyle=\bfseries\color{keyword},
  commentstyle=\itshape\color{comment},
  stringstyle=\color{string},
  numberstyle=\tiny\color{gray},
  numbers=left,
  numbersep=6pt,
  showstringspaces=false,
  breaklines=true,
  breakatwhitespace=true,
  tabsize=4,
  captionpos=b,
  morekeywords={True,False,None,np,plt,pd,os,sys,mp,datetime},
}
\lstset{style=python}

% ── Formatação de títulos ─────────────────────────────────────────────────────
\titleformat{\chapter}[display]
  {\normalfont\huge\bfseries\color{cinzatitulo}}
  {\chaptertitlename\ \thechapter}{18pt}{\Huge}
\titleformat{\section}
  {\normalfont\Large\bfseries\color{cinzatitulo}}
  {\thesection}{1em}{}
\titleformat{\subsection}
  {\normalfont\large\bfseries\color{cinzatitulo!80}}
  {\thesubsection}{1em}{}
\titleformat{\subsubsection}
  {\normalfont\normalsize\bfseries}
  {\thesubsubsection}{1em}{}

% ── Cabeçalho e rodapé ───────────────────────────────────────────────────────
\pagestyle{fancy}
\fancyhf{}
\fancyhead[L]{\small\textit{Natalia Time — Documentação Técnica}}
\fancyhead[R]{\small\textit{\leftmark}}
\fancyfoot[C]{\thepage}
\renewcommand{\headrulewidth}{0.4pt}

% ── Espaçamento ───────────────────────────────────────────────────────────────
\setlength{\parskip}{6pt}
\setlength{\parindent}{0pt}

% ── Comandos auxiliares ──────────────────────────────────────────────────────
\newcommand{\code}[1]{\texttt{#1}}
\newcommand{\param}[1]{\texttt{\bfseries #1}}
\newcommand{\arquivo}[1]{\texttt{\itshape #1}}

% =============================================================================
\begin{document}

\begin{titlepage}
  \centering
  \vspace*{2cm}
  {\Huge\bfseries\color{cinzatitulo} Natalia Time}\par
  \vspace{0.5cm}
  {\large\color{azul} Documentação Técnica Completa — v1.1.0}\par
  \vspace{1cm}
  \rule{0.8\textwidth}{1pt}\par
  \vspace{0.5cm}
  {\large Ferramenta de Análise de Espectroscopia de Tempo de Voo (DETOF)\\
  para Fragmentação por Impacto de Elétrons}\par
  \vspace{1.5cm}
  {\normalsize
    Cobre os três componentes principais do sistema:\par
    \vspace{0.3cm}
    \texttt{analise\_natalia\_time.py} — Motor de análise e otimização\par
    \texttt{NT\_app\_ctk.py} — Interface gráfica (customtkinter)\par
    \texttt{Migrador NT/migrador.py} — Conversor de formatos\par
  }
  \vspace{2cm}
  \rule{0.8\textwidth}{0.4pt}\par
  \vspace{0.5cm}
  {\small Documento gerado em março de 2026.\\
  Para uso interno no Laboratório de Física Molecular e Espectroscopia.}
\end{titlepage}

\tableofcontents
\newpage

% =============================================================================
\chapter{Contexto Físico}
\label{cap:fisica}
% =============================================================================

\section{Espectroscopia de Tempo de Voo (DETOF)}

A técnica DETOF (\textit{Drift-time Electrostatic Time-of-Flight}) mede o
tempo que íons fragmentos levam para percorrer uma distância conhecida após
serem produzidos em uma câmara de ionização por impacto de elétrons.
A variável medida é, portanto, o tempo $t$ (em nanosegundos) e o sinal
correspondente $u(t)$, proporcional à quantidade de íons que chegam ao
detector em cada instante.

O experimento tipicamente produz uma curva $u/u_0$ versus $t$ em escala
logarítmica, onde $u_0$ é o sinal no tempo de referência $t_0 =$ \texttt{I7\_NORM}
(padrão: 500~ns). A normalização é essencial para comparar experimentos
realizados em diferentes condições de corrente de feixe.

\subsection{Fragmentação por Impacto de Elétrons}

Quando uma molécula neutra (por exemplo, O$_2$) é bombardeada por elétrons
energéticos (tipicamente 100--300~eV), ocorrem ionização e fragmentação:
\[
  e^- + \text{O}_2 \longrightarrow \text{O}^+ + \text{O} + 2e^-
\]
O fragmento iônico O$^+$ é produzido com uma distribuição de energias
cinéticas que reflete o mecanismo de fragmentação: parte dos fragmentos
carrega a energia térmica do gás de fundo (distribuição de Maxwell-Boltzmann),
enquanto outros recebem energia cinética adicional proveniente da quebra de
ligações, resultando em distribuições exponenciais e gaussianas.

\subsection{Por Que Usar uma Soma Ponderada de Distribuições?}

A distribuição total de velocidades do fragmento não é descrita por uma
única função analítica. Diferentes canais de fragmentação produzem fragmentos
com características energéticas distintas:

\begin{itemize}[leftmargin=1.5em]
  \item \textbf{Maxwell-Boltzmann (MB):} fragmentos em equilíbrio térmico com
    o gás ambiente — presença inevitável em qualquer experimento.
  \item \textbf{Exponencial (EXP):} fragmentos com distribuição exponencial de
    energia cinética — canal dissociativo "suave".
  \item \textbf{Gaussianas (G1, G2, G3):} fragmentos com energia bem definida,
    emitidos por estados específicos da molécula-mãe.
\end{itemize}

A modelagem por soma ponderada permite identificar e quantificar a contribuição
relativa de cada canal, mesmo quando eles se sobrepõem no espectro de tempo.

\section{O Modelo Físico}

O sinal teórico $S(t)$ é expresso como:

\begin{equation}
  \boxed{S(t) = p_{\text{mb}}\cdot\text{MB}(t) + p_{\text{exp}}\cdot\text{EXP}(t)
               + p_{g1}\cdot G_1(t) + p_{g2}\cdot G_2(t) + p_{g3}\cdot G_3(t)}
  \label{eq:modelo}
\end{equation}

com a restrição de normalização:
\begin{equation}
  p_{\text{mb}} + p_{\text{exp}} + p_{g1} + p_{g2} + p_{g3} = 1
  \label{eq:norma}
\end{equation}

Cada parcela $X(t)$ representa o sinal esperado de uma curva isolada,
normalizado de forma independente para que sua amplitude no ponto de
referência $t_0$ seja unitária. O código implementa isso via a função
\code{normalizar\_amplitudes()}, que replica exatamente o procedimento
VBA \textit{NormalizarProcessoCientifico} da planilha original.

\subsection{Componente Maxwell-Boltzmann}

A distribuição de velocidades de Maxwell-Boltzmann para fragmentos em
equilíbrio térmico com o gás à temperatura $T$ é:
\[
  f_{\text{MB}}(v) \propto v \cdot \exp\!\left(-\beta_{\text{mb}}^2 v^2\right)
\]

O parâmetro $\beta_{\text{mb}}$ está relacionado à temperatura por:
\begin{equation}
  \beta_{\text{mb}} = \sqrt{\frac{4{,}04 \times 10^5 \cdot (300/T_{\text{gas}}) \cdot m_{\text{mol}}}{2}}
  \label{eq:betamb}
\end{equation}

onde $T_{\text{gas}}$ é a temperatura do gás em kelvin e $m_{\text{mol}}$ é
a massa molar da molécula-mãe. O parâmetro $\beta_{\text{mb}}$ é lido
diretamente da célula AA4 da planilha Excel (ou do \texttt{parametros.toml})
e é \textbf{fixo durante toda a otimização} — ele é determinado pela física
do experimento, não pelo ajuste.

No código, a distribuição MB pura é calculada como:
\begin{lstlisting}[caption={Distribuição MB na grade de velocidades}]
y_mb = H9_VEC * np.exp(-((fixos["beta_mb"] * H9_VEC)**2))
\end{lstlisting}

onde \code{H9\_VEC} é o vetor de velocidades discretas na grade de
Gauss-Hermite.

\subsection{Componente Exponencial}

A distribuição exponencial modela fragmentos que recebem energia cinética
adicional por um canal dissociativo com distribuição de energia exponencial:
\[
  f_{\text{exp}}(v) \propto v \cdot \exp\!\left(-\beta_{\text{exp}}^2 v^2\right)
  \cdot \text{GH}_{\text{exp}}(\beta_{\text{exp}})
\]

O kernel de Gauss-Hermite \code{gh\_exp} integra numericamente:
\[
  \text{GH}_{\text{exp}}(\beta^2) = \sum_{k=1}^{4} w_k \cdot
    \frac{e^{-x_k^2} \cdot x_k}{\sqrt{x_k^4 + x_k^2 \beta^2 v^2}}
\]

O parâmetro $\beta_{\text{exp}}$ é \textbf{variável} e otimizado durante
a busca, no intervalo configurável (padrão: 50 a 5000).

\subsection{Componentes Gaussianas}

Cada gaussiana $G_i(t)$ corresponde a fragmentos emitidos com energia cinética
bem definida $E_i$ (em eV). A velocidade central correspondente é:
\begin{equation}
  v_0 = \frac{\sqrt{E_i / m_{\text{frag}}}}{72{,}284}
  \label{eq:v0}
\end{equation}

A distribuição no espaço de velocidades é:
\[
  f_{G}(v) \propto \frac{v}{2} \cdot \exp\!\left(-\beta_G^4 (v^2 - v_0^2)^2\right)
  \cdot \text{GH}_{\text{gauss}}(\beta_G, v_0)
\]

As energias $E_{g1}$, $E_{g2}$, $E_{g3}$ (variáveis \code{e\_g3}, \code{e\_g4},
\code{e\_g5} no código) são tipicamente lidas da planilha mas podem ser
sobrescritas pela configuração. Cada gaussiana tem seu próprio parâmetro
de largura $\beta_{Gi}$ otimizado.

\subsection{A Restrição de Normalização dos Pesos}

O código normaliza os pesos automaticamente:
\begin{lstlisting}[caption={Normalização dos pesos em calcular\_espectro()}]
soma_p = p_mb + p_exp + p_g3 + p_g4 + p_g5
pn_mb  = p_mb  / soma_p
pn_exp = p_exp / soma_p
# etc.
\end{lstlisting}

Isso garante que a equação~\ref{eq:norma} seja sempre satisfeita, independente
dos valores brutos sorteados. Na busca aleatória, os pesos são amostrados
de uma distribuição de Dirichlet, que por construção produz vetores que
somam~1.

\section{Métricas de Qualidade: RMSE e R\texorpdfstring{$^2$}{2}}

\subsection{Coeficiente de Determinação R\texorpdfstring{$^2$}{2}}

O $R^2$ replica exatamente a fórmula Excel usada na planilha de referência:
\begin{equation}
  R^2 = 1 - \frac{\sum_i (S_{\text{teo},i} - S_{\text{exp},i})^2}
                 {\sum_i (S_{\text{exp},i} - \bar{S}_{\text{exp}})^2}
  \label{eq:r2}
\end{equation}

Um resultado é considerado válido quando $R^2 \geq$ \code{SCORE\_MINIMO}
(padrão: 0,99). Valores de $R^2$ próximos de 1,0 indicam excelente concordância.

\subsection{RMSE Relativo}

O RMSE implementado não é o erro quadrático médio simples. Ele replica o
algoritmo VBA \textit{CalcularDesvioResiduos}, que usa resíduos relativos
com subtração de tendência linear:

\begin{enumerate}
  \item Calcula o resíduo absoluto: $D_i = S_{\text{teo},i} - S_{\text{exp},i}$
  \item Ajusta uma reta sobre $D_i$ vs. $t_i$: $D_i \approx m \cdot t_i + b$
  \item Calcula o resíduo relativo: $E_i = D_i / S_{\text{exp},i}$
  \item Retorna:
  \[
    \text{RMSE} = \sqrt{\frac{1}{N}\sum_i (E_i - (m \cdot t_i + b))^2}
  \]
\end{enumerate}

Essa definição é mais sensível a desvios sistemáticos (tendências no resíduo)
do que um RMSE absoluto simples, e é por isso que a função objetivo do Nelder-Mead
minimiza o RMSE e não o $1 - R^2$.

\section{Parâmetros Físicos do Experimento}

\begin{table}[h!]
  \centering
  \caption{Parâmetros físicos lidos da planilha ou do \texttt{parametros.toml}}
  \begin{tabular}{llp{7cm}}
    \toprule
    \textbf{Variável} & \textbf{Célula Excel} & \textbf{Descrição} \\
    \midrule
    \code{beta\_mb} & AA4 (linha 4, col 27) & Parâmetro de largura da distribuição MB; determina a temperatura efetiva do gás \\
    \code{m\_frag}  & B14 (linha 14, col 2) & Massa do fragmento iônico (u.a.m.) \\
    \code{m\_mol}   & B13 (linha 13, col 2) & Massa da molécula-mãe (u.a.m.) \\
    \code{DE}       & B5  (linha 5, col 2)  & Tensão de extração do espectrômetro (eV) \\
    \code{D}        & B9  (linha 9, col 2)  & Diâmetro do tubo de voo (m) \\
    \code{LL}       & B11 (linha 11, col 2) & Comprimento efetivo do cilindro de ionização (m) \\
    \code{I7\_NORM} & I7  (linha 7, col 9)  & Tempo de referência para normalização das amplitudes (ns) \\
    \code{e\_g3}    & AC6 (linha 6, col 29) & Energia da Gaussiana 1 (eV) \\
    \code{e\_g4}    & AD6 (linha 6, col 30) & Energia da Gaussiana 2 (eV) \\
    \code{e\_g5}    & AE6 (linha 6, col 31) & Energia da Gaussiana 3 (eV) \\
    \bottomrule
  \end{tabular}
\end{table}

\subsection{Valores de Fallback}

Quando uma célula da planilha está vazia ou inacessível, o código utiliza
\code{FIXOS\_FALLBACK}, uma constante definida na Seção~4 de
\arquivo{analise\_natalia\_time.py}. Esses valores correspondem ao experimento
de referência O$_2$ a 300~eV:

\begin{lstlisting}[caption={Dicionário FIXOS\_FALLBACK}]
FIXOS_FALLBACK = {
    "beta_mb": 2572.6309556813,
    "e_g3":    0.7,
    "e_g4":    3.0,
    "e_g5":    2.0,
    "m_frag":  16.0,
    "m_mol":   32.0,
    "DE":      350.0,
    "D":       6.8,
    "LL":      6.8,
    "I7_NORM": 500.0,
}
\end{lstlisting}

\section{A Grade BM e a Interpolação}

O sinal teórico $S(t)$ envolve uma integral sobre todas as velocidades
possíveis dos fragmentos. Para calcular essa integral numericamente e
depois comparar com os pontos experimentais (que podem ter espaçamento
irregular), o código usa a seguinte estratégia:

\begin{enumerate}
  \item \textbf{Construção da grade BM} (\code{GRADE\_BM}): um vetor de tempos
    com três regimes de resolução que replicam a planilha Excel original:
    \begin{itemize}
      \item 510–2960~ns: passo de 50~ns
      \item 3060–9960~ns: passo de 100~ns
      \item 10160~ns em diante: passo de 200~ns
    \end{itemize}
    Tipicamente ~175 pontos, cobrindo todo o range experimental.

  \item \textbf{Cálculo vetorial} (\code{calcular\_espectro}): para cada ponto
    da grade, o sinal é calculado usando operações vetorizadas com
    \texttt{numpy.einsum}, que substitui um loop Python de 175 iterações
    por uma operação de álgebra linear muito mais rápida.

  \item \textbf{Interpolação linear} (\code{numpy.interp}): o sinal na grade
    BM é interpolado nos tempos experimentais reais, que podem ter
    espaçamento diferente.
\end{enumerate}

O motivo de usar a grade ao invés de calcular diretamente nos tempos
experimentais é que os filtros de aceitação do detector $P_9(t)$ e $U_9(t)$
são pré-calculados para toda a grade e armazenados em \code{P9\_GRADE} e
\code{U9\_GRADE} (matrizes de dimensão $175 \times 1492$), evitando
recálculos custosos a cada iteração da otimização.


% =============================================================================
\chapter{Pipeline de Otimização}
\label{cap:otimizacao}
% =============================================================================

\section{Visão Geral}

O pipeline de otimização tem duas fases complementares:

\begin{description}[leftmargin=2.5em, style=nextline]
  \item[\textbf{Fase 1 — Busca global aleatória}]
    Explora o espaço de parâmetros de forma ampla, usando paralelismo
    via \code{multiprocessing.Pool}. Para cada ``rodada'', sorteia um
    conjunto de $\beta$s aleatórios e testa múltiplos conjuntos de pesos.
    Guarda todos os resultados com $R^2 \geq$ \code{SCORE\_MINIMO}.

  \item[\textbf{Fase 2 — Refinamento local com Nelder-Mead}]
    Seleciona os melhores candidatos da busca aleatória e os refina com
    o método de Nelder-Mead (simplex), que é um otimizador local sem
    gradiente — adequado para funções com descontinuidades numéricas.
\end{description}

Esse design é uma estratégia clássica de \textit{global search + local
refinement}: a busca aleatória evita ficar preso em mínimos locais,
enquanto o Nelder-Mead extrai o máximo de qualidade dos candidatos
promissores.

\section{Fase 1: Busca Global Aleatória}

\subsection{Amostragem de Betas}

Para cada rodada $i$ (de 1 a \code{ITERACOES\_BETA}), os quatro betas
variáveis são sorteados independentemente de distribuições uniformes:
\[
  \beta_k \sim \mathcal{U}(lo_k,\; hi_k), \quad k \in \{\text{exp}, g1, g2, g3\}
\]

Os intervalos $[lo_k, hi_k]$ são configuráveis via \code{INTERVALOS\_BETA}.
O beta da MB ($\beta_{\text{mb}}$) é fixo e lido da planilha.

\subsection{Amostragem de Pesos com Distribuição de Dirichlet}

Para cada conjunto de $\beta$s, são testadas \code{ITERACOES\_PESOS}
combinações de pesos. Os pesos são amostrados de uma distribuição de
Dirichlet com parâmetro de concentração uniforme:
\[
  (p_1, p_2, \ldots, p_K) \sim \text{Dirichlet}(\mathbf{1}_K)
\]

A distribuição de Dirichlet garante que os pesos sejam positivos e somem 1
por construção, sem necessidade de normalização posterior (embora o código
a execute assim mesmo por segurança). Ela também produz amostras uniformes
sobre o simplex — crucial para explorar o espaço de pesos sem viés.

\begin{lstlisting}[caption={Amostragem de pesos via Dirichlet}]
curvas_ativas = [k for k, v in CURVAS_PESO_ATIVAS.items() if v]
n_curvas      = len(curvas_ativas)
for _ in range(ITERACOES_PESOS):
    dir_s = rng.dirichlet(np.ones(n_curvas))
    pesos = {k: 0.0 for k in nomes_pesos}
    for idx, nome in enumerate(curvas_ativas):
        pesos[nome] = float(dir_s[idx])
\end{lstlisting}

\subsection{Critério de Aceitação}

Um resultado é marcado como \textbf{VÁLIDO} se:
\[
  R^2 \geq \text{SCORE\_MINIMO} \quad (\text{padrão: } 0{,}99)
\]

O cálculo de RMSE é feito \emph{apenas} para resultados válidos, pois é
mais caro computacionalmente. Para resultados inválidos, apenas $R^2$
é calculado (muito mais rápido).

\subsection{Estrutura de Paralelismo com multiprocessing.Pool}

O paralelismo é implementado via \code{multiprocessing.Pool} com o
padrão \texttt{imap\_unordered}:

\begin{lstlisting}[caption={Pool multiprocessing com imap\_unordered}]
pool = mp.Pool(processes=n_proc,
               initializer=_worker_init,
               initargs=(fixos, tempos_exp, sinal_exp))
for reg in pool.imap_unordered(_processar_rodada,
                               args_list, chunksize=chunk):
    historico.append(...)
\end{lstlisting}

Cada processo worker é inicializado com \code{\_worker\_init}, que
reconstrói as grades \code{P9\_GRADE} e \code{U9\_GRADE} localmente
(necessário porque objetos NumPy não são passados eficientemente por
\texttt{pickle} entre processos). Os workers são independentes —
cada um processa uma rodada completa.

A função \code{\_processar\_rodada} é a unidade de trabalho: recebe
o índice da rodada e uma semente aleatória, e retorna o melhor
resultado encontrado nessa rodada.

\subsection{Mecanismo de Parada Suave (PARAR.txt)}

Para permitir interrupção sem perda de dados, o código verifica
periodicamente a existência de um arquivo \code{PARAR.txt} na pasta
de resultados:

\begin{lstlisting}[caption={Verificação de arquivo PARAR.txt}]
arquivo_parar = os.path.join(pasta_parar, ARQUIVO_PARAR)
if os.path.exists(arquivo_parar):
    print(f"[{ARQUIVO_PARAR} detectado] Encerrando...")
    interrompido = True
    break
\end{lstlisting}

A criação do arquivo pode ser feita pelo botão ``Parar'' da GUI
ou manualmente. O código para ao final da rodada atual, salvando
todos os resultados coletados.

\section{Fase 2: Refinamento Local com Nelder-Mead}

\subsection{Seleção de Candidatos por Clustering Hierárquico}

Após a busca aleatória, os resultados válidos são agrupados no espaço
normalizado dos betas (4 dimensões, cada uma em $[0, 1]$). O objetivo
é selecionar candidatos \emph{diversos} para o Nelder-Mead, não apenas
os $N$ melhores por RMSE que podem estar todos no mesmo mínimo local.

O algoritmo usa clustering hierárquico aglomerativo com linkagem completa
(\texttt{scipy.cluster.hierarchy}):
\[
  d_{\text{normalizada}} = \frac{\beta_k - lo_k}{hi_k - lo_k}
\]

O corte é definido por \code{NM\_CLUSTER\_LIMIAR} (padrão: 0,25) aplicado
à distância euclidiana. Do melhor elemento de cada cluster (menor RMSE),
é selecionado um candidato.

\begin{figure}[h!]
\centering
\begin{minipage}{0.85\textwidth}
\begin{description}[leftmargin=2em, style=nextline]
  \item[\textbf{Exemplo com 3 clusters:}]
    Suponha que a busca encontrou 50 resultados válidos. Dois deles
    têm $\beta_{\text{exp}} \approx 300$ e um tem $\beta_{\text{exp}} \approx 800$.
    O clustering os separa em 2 grupos e seleciona 1 candidato de cada,
    garantindo que o Nelder-Mead explore ambas as regiões do espaço.
\end{description}
\end{minipage}
\end{figure}

\subsection{Parametrização do Nelder-Mead}

O Nelder-Mead otimiza simultaneamente:
\begin{itemize}
  \item Os 4 betas variáveis: $\beta_{\text{exp}}, \beta_{g1}, \beta_{g2}, \beta_{g3}$
  \item Os pesos: parametrizados como \textbf{logits} (softmax), garantindo
    que somem 1 sem restrições explícitas
  \item Opcionalmente: as energias das gaussianas (quando \code{ENERGIA\_NM\_LIVRE = True})
\end{itemize}

A parametrização por softmax é uma escolha de design importante:
permite ao otimizador trabalhar em espaço irrestrito ($\mathbb{R}^K$)
enquanto os pesos resultantes ficam automaticamente em $[0,1]$ e somam~1.

\begin{lstlisting}[caption={Softmax com piso de peso mínimo}]
def softmax_com_piso(w):
    e = np.exp(w - np.max(w))
    probs = e / e.sum()
    for idx, nome in enumerate(curvas_ativas):
        pmin = PESO_MIN.get(nome)
        if pmin is not None and probs[idx] < pmin:
            probs[idx] = pmin
    return probs / probs.sum()
\end{lstlisting}

\subsection{Reinícios Automáticos}

Se o Nelder-Mead não converge em \code{NM\_MAX\_ITER} iterações, o código
o reinicia com perturbação do ponto ótimo encontrado:
\[
  x_{\text{novo}} = x_{\text{melhor}} + \epsilon \cdot \delta
\]
onde $\delta$ é amostrado de $\mathcal{U}(-1, 1)$ e $\epsilon = $
\code{NM\_PERTURB\_ESCALA}~$ \times$ escala do parâmetro (padrão: $\pm 10\%$).

\subsection{Nomenclatura das Rodadas no Output}

\begin{table}[h!]
  \centering
  \caption{Convenção de nomenclatura das rodadas}
  \begin{tabular}{lp{8cm}}
    \toprule
    \textbf{Prefixo} & \textbf{Significado} \\
    \midrule
    \texttt{1, 2, 3, ...} & Rodadas numéricas da busca aleatória \\
    \texttt{NM1, NM2, ...} & Rodadas do refinamento Nelder-Mead (candidato 1, 2, ...) \\
    \texttt{F2-N} & Rodadas da Fase 2 (busca em duas fases) \\
    \texttt{A-N} & Rodadas da busca adaptada \\
    \texttt{VAL} & Rodada do modo de validação manual \\
    \bottomrule
  \end{tabular}
\end{table}

\section{Arquivos de Saída}

Ao final de cada execução, o código cria uma subpasta em
\texttt{Resultados/<nome\_experimento>\_<data>\_<hora>/} contendo:

\begin{description}[leftmargin=2.5em, style=nextline]
  \item[\texttt{<nome>\_relatorio.txt}]
    Relatório de texto plano com metadados, configuração usada e tabela
    dos melhores resultados.

  \item[\texttt{<nome>\_Relatorio<N>\_<stamp>.xlsx}]
    Relatório Excel com seções de metadados (início, duração, válidos),
    configuração de NM, intervalos de busca, melhor resultado global e
    tabela completa de todas as rodadas ordenadas por $R^2$.
    Células verdes = válido; vermelhas = abaixo do corte.

  \item[\texttt{graficos\_rodadas/rod<NNNN>\_*.png}]
    Gráfico para cada rodada válida (limitado a \code{SALVAR\_GRAFICOS\_TOP},
    padrão: 20). Cada gráfico tem dois painéis: espectro log-log com as
    curvas componentes e resíduos relativos com reta de tendência.
\end{description}


% =============================================================================
\chapter{Planilhas Excel de Entrada e Saída}
\label{cap:excel}
% =============================================================================

\section{Estrutura da Planilha de Entrada (.xlsm/.xlsx)}

\subsection{Dados Experimentais}

Os dados experimentais ocupam as colunas A e B da primeira aba da planilha:
\begin{itemize}
  \item \textbf{Coluna A}: tempo em nanossegundos
  \item \textbf{Coluna B}: sinal normalizado $u/u_0$
\end{itemize}

O código detecta automaticamente a linha de início procurando a palavra
``tempo'' (insensível a maiúsculas/minúsculas) na coluna A:

\begin{lstlisting}[caption={Detecção automática da linha de início}]
for i, row in enumerate(ws.iter_rows(...), start=1):
    val = row[0]
    if val is not None and isinstance(val, str) \
       and "tempo" in val.lower():
        return i + 1  # linha seguinte = início dos dados
\end{lstlisting}

A leitura é feita em \textbf{passagem única} (\texttt{read\_only=True}),
necessária para não carregar a planilha inteira na memória.

\subsection{Parâmetros Fixos: Mapeamento de Células}

\begin{table}[h!]
  \centering
  \caption{Células fixas lidas da planilha (índices base-0 no código)}
  \begin{tabular}{llll}
    \toprule
    \textbf{Variável} & \textbf{Índice (row, col)} & \textbf{Célula Excel} & \textbf{Valor típico} \\
    \midrule
    \code{beta\_mb} & (3, 26) & AA4  & 2572,63 \\
    \code{m\_frag}  & (13, 1) & B14  & 16,0 \\
    \code{m\_mol}   & (12, 1) & B13  & 32,0 \\
    \code{DE}       & (4, 1)  & B5   & 350,0 eV \\
    \code{D}        & (8, 1)  & B9   & 6,8 m \\
    \code{LL}       & (10, 1) & B11  & 6,8 m \\
    \code{I7\_NORM} & (6, 8)  & I7   & 500,0 ns \\
    \code{e\_g3}    & (5, 28) & AC6  & 0,7 eV \\
    \code{e\_g4}    & (5, 29) & AD6  & 3,0 eV \\
    \code{e\_g5}    & (5, 30) & AE6  & 2,0 eV \\
    \bottomrule
  \end{tabular}
\end{table}

Os índices \texttt{(row, col)} usam base-0; a célula Excel é obtida somando
1 a cada índice no código \code{ws.cell(row=r+1, column=c+1)}.

\section{Formato CSV + TOML}

O formato alternativo (criado pelo Migrador) consiste em dois arquivos
na mesma pasta:

\subsection{dados.csv}

Arquivo de texto com duas colunas separadas por vírgula:
\begin{verbatim}
tempo,sinal
510.0,1.0
560.0,0.985348
610.0,0.950200
...
\end{verbatim}

O código aceita tanto vírgula quanto ponto-e-vírgula como separador.
Linhas não numéricas (incluindo o cabeçalho) são ignoradas silenciosamente.

\subsection{parametros.toml}

Arquivo TOML com os parâmetros físicos. Chaves obrigatórias:
\code{beta\_mb}, \code{m\_frag}, \code{m\_mol}. Chaves opcionais:
\code{e\_g3}, \code{e\_g4}, \code{e\_g5}, \code{DE}, \code{D},
\code{LL}, \code{I7\_NORM}.

\begin{lstlisting}[language={},caption={Exemplo de parametros.toml}]
# Parametros fisicos do experimento
beta_mb = 2572.630956
m_frag  = 16.0
m_mol   = 32.0

# Parametros geometricos do espectrometro DETOF
DE      = 350.0
D       = 6.8
LL      = 6.8
I7_NORM = 500.0

# Energias das gaussianas (eV)
# e_g3 = 0.7
# e_g4 = 3.0
# e_g5 = 2.0
\end{lstlisting}

As energias ficam comentadas por padrão; são descomentadas e preenchidas
pelo botão ``Gravar'' da GUI após o usuário inserir os valores.

\section{Estrutura do Relatório Excel de Saída}

O arquivo \texttt{<nome>\_Relatorio<N>\_<stamp>.xlsx} tem uma única aba
com a seguinte estrutura linear (linha a linha):

\begin{enumerate}
  \item \textbf{Seção METADADOS DA OTIMIZAÇÃO}: data/hora de início e fim,
    duração, iterações configuradas, total de cálculos, válidos, alta precisão.

  \item \textbf{Seção REFINAMENTO NELDER-MEAD}: se ativo — número de
    candidatos, máximo de iterações, reinícios, tolerância, rodadas NM
    e quantas foram válidas.

  \item \textbf{Seção INTERVALOS DE BUSCA (Beta)}: para cada beta, a faixa
    $[\text{min}, \text{max}]$ usada. Seguida de seção ENERGIAS DAS GAUSSIANAS.

  \item \textbf{Seção MELHOR RESULTADO GLOBAL}: todos os parâmetros do melhor
    ajuste encontrado (rodada, R², RMSE, equação da reta de resíduos, betas,
    energias, pesos).

  \item \textbf{Tabela completa}: cabeçalho com fundo azul-escuro; cada rodada
    em uma linha. Rodadas válidas em verde claro (\texttt{D5F5E3}); inválidas
    em rosa (\texttt{FADBD8}).
\end{enumerate}

\subsection{Colunas da Tabela Completa}

\begin{table}[h!]
  \centering
  \caption{Colunas da tabela de resultados no relatório Excel}
  \begin{tabular}{ll}
    \toprule
    \textbf{Coluna} & \textbf{Conteúdo} \\
    \midrule
    Rodada         & Identificador (número, NM1, NM2, etc.) \\
    Status         & \texttt{VÁLIDO} ou \texttt{Abaixo do Corte} \\
    R²             & Coeficiente de determinação \\
    RMSE           & Erro quadrático médio relativo \\
    Equação da Reta & $y = m \cdot x + b$ dos resíduos \\
    Beta Exp       & $\beta_{\text{exp}}$ otimizado \\
    Beta G1        & $\beta_{g1}$ otimizado \\
    Beta G2        & $\beta_{g2}$ otimizado \\
    Beta G3        & $\beta_{g3}$ otimizado \\
    Energia G1 (eV) & $E_{g1}$ usada \\
    Energia G2 (eV) & $E_{g2}$ usada \\
    Energia G3 (eV) & $E_{g3}$ usada \\
    Peso MB        & $p_{\text{mb}}$ normalizado \\
    Peso Exp       & $p_{\text{exp}}$ normalizado \\
    Peso G1        & $p_{g1}$ normalizado \\
    Peso G2        & $p_{g2}$ normalizado \\
    Peso G3        & $p_{g3}$ normalizado \\
    \bottomrule
  \end{tabular}
\end{table}


% =============================================================================
\chapter{analise\_natalia\_time.py — Referência Detalhada}
\label{cap:analise}
% =============================================================================

\section{Visão Geral e Estrutura em Seções}

O arquivo \arquivo{analise\_natalia\_time.py} é o motor de análise e pode
ser executado diretamente (\code{python analise\_natalia\_time.py}) ou
como subprocesso gerado pela GUI. Está organizado em 13 seções bem delimitadas:

\begin{table}[h!]
  \centering
  \caption{Estrutura de seções de analise\_natalia\_time.py}
  \begin{tabular}{rp{11cm}}
    \toprule
    \textbf{Seção} & \textbf{Conteúdo} \\
    \midrule
    1  & Importações e utilitários (função \code{br()} para formatação brasileira) \\
    2  & \textbf{Configuração} — todos os parâmetros injetados pela GUI \\
    3  & Constantes físicas e geométricas; grade de Gauss-Hermite \\
    4  & Leitura do Excel e do CSV+TOML \\
    5  & Kernels de Gauss-Hermite (\code{gh\_exp}, \code{gh\_gauss}) \\
    6  & Filtros de aceitação do detector (\code{calcular\_filtros}) \\
    7  & Normalização das amplitudes (\code{normalizar\_amplitudes}) \\
    8  & Integração do sinal teórico (\code{calcular\_espectro}) \\
    9  & Métricas R² e RMSE \\
    10 & Busca aleatória global (worker, pool, parada) \\
    11 & Gráficos e viewer interativo \\
    12 & Relatório Excel (\code{escrever\_relatorio\_excel}) \\
    13 & Execução principal (\code{\_rodar\_analise}, \code{main}) \\
    \bottomrule
  \end{tabular}
\end{table}

A Seção~2 é especial: a GUI a \textbf{substitui completamente} ao gerar o
script temporário, injetando os valores configurados pelo usuário.

\section{Constantes Físicas Globais}

\begin{description}[leftmargin=2.5em, style=nextline]
  \item[\texttt{FATOR\_ENERGIA = 7837.5}]
    Converte velocidade em energia: $E = \text{FATOR\_ENERGIA} \cdot m \cdot v^2$.
  \item[\texttt{FATOR\_VELOCIDADE = 16.8}]
    Fator de conversão de velocidade para cálculo de tempos de chegada
    ao detector após a região de aceleração.
  \item[\texttt{H6 = 120000}]
    Número total de pontos na grade de velocidades (resolução da grade).
    \texttt{H6\_INV = 1/H6} é o passo de velocidade.
  \item[\texttt{N\_VEL = 1492}]
    Número de pontos de velocidade efetivamente usados na integração
    numérica de Gauss-Hermite.
  \item[\texttt{FMB = 1.7}]
    Fator multiplicativo específico da curva MB (replica constante da planilha).
  \item[\texttt{X\_GH, W\_GH}]
    Nós e pesos de Gauss-Hermite de 4 pontos — escolhidos para integrar
    exatamente polinômios de grau $\leq 7$, suficiente para as distribuições
    em questão.
\end{description}

\section{Função inicializar\_grade()}

\begin{lstlisting}[caption={Assinatura de inicializar\_grade()}]
def inicializar_grade(fixos: dict, t_min_exp: float = 510.0,
                      t_max_exp: float = 19960.0,
                      verbose: bool = True):
\end{lstlisting}

\textbf{Responsabilidade:} Inicializa todas as variáveis globais dependentes
dos parâmetros físicos. Deve ser chamada \emph{antes} de qualquer cálculo
de espectro.

\textbf{O que faz:}
\begin{enumerate}
  \item Calcula \code{L9\_VEC} e \code{M9\_VEC}: energias do fragmento e da
    molécula-mãe na grade de velocidades ($E = \text{FATOR\_ENERGIA} \cdot m \cdot v^2$).
  \item Atualiza os globais geométricos (\code{D}, \code{DD}, \code{LL},
    \code{DE}, \code{I7\_NORM}, \code{N0}) com os valores lidos da planilha.
  \item Constrói \code{GRADE\_BM}: garante cobertura do range experimental
    com resolução variável (50/100/200~ns).
  \item Pré-calcula \code{P9\_GRADE} e \code{U9\_GRADE} (matrizes
    $175 \times 1492$) com todos os filtros para cada ponto da grade.
    Isso é a operação mais cara da inicialização.
  \item Calcula \code{P9\_NORM} e \code{U9\_NORM} para o ponto de
    normalização \code{I7\_NORM}.
\end{enumerate}

\textbf{Por que pré-calcular os filtros?} Os filtros $P_9(t)$ e $U_9(t)$
dependem do tempo $t$ mas não dos betas — eles modelam a geometria do
espectrômetro. Pré-calculá-los uma única vez economiza milhões de chamadas
a \code{calcular\_filtros()} durante a otimização.

\section{Função calcular\_espectro()}

\begin{lstlisting}[caption={Assinatura de calcular\_espectro()}]
def calcular_espectro(tempos_exp, beta_exp, beta_g3, beta_g4, beta_g5,
                      p_mb, p_exp, p_g3, p_g4, p_g5, fixos, amps=None):
\end{lstlisting}

Esta é a função computacionalmente mais crítica — é chamada
$\text{ITERACOES\_BETA} \times \text{ITERACOES\_PESOS}$ vezes por execução.

\textbf{Estratégia de otimização:}
\begin{enumerate}
  \item Normaliza os pesos por sua soma.
  \item Reutiliza os vetores $y$ das distribuições puras calculados em
    \code{normalizar\_amplitudes()} — evita recalcular os kernels de
    Gauss-Hermite, que são as operações mais custosas.
  \item Usa \texttt{numpy.einsum} para a integral vetorial, substituindo
    um loop Python de 175 iterações:
    \begin{lstlisting}[caption={Integração vetorial com einsum}]
# P9_GRADE shape: (175, 1492)
# y_e * dv shape: (1492,)
sinal_grade += pn_exp * amps["ak3"] * \
    np.einsum('ij,j->i', P9_GRADE, y_e * dv)
    \end{lstlisting}
  \item Interpola o resultado na grade BM nos tempos experimentais reais
    com \texttt{numpy.interp}.
\end{enumerate}

\section{Função funcao\_objetivo() (dentro do NM)}

A função objetivo usada pelo Nelder-Mead é:
\begin{lstlisting}[caption={Função objetivo do Nelder-Mead}]
def objetivo(x):
    betas, energias, pesos = params_de_x(x)
    fixos_nm = {**fixos, **energias}
    amps = normalizar_amplitudes(fixos_nm, ...)
    sinal_teo = calcular_espectro(tempos_exp, ..., fixos_nm, amps)
    rmse, _, _ = calcular_rmse(tempos_exp, sinal_exp, sinal_teo)
    return rmse
\end{lstlisting}

O motivo de minimizar RMSE (e não $1 - R^2$) é que o RMSE é mais sensível
a desvios finos próximos do ótimo. O $R^2$ é usado para o filtro de validade
porque é mais intuitivo para o usuário.

\section{Função busca\_aleatoria\_worker()}

A função \code{\_processar\_rodada(args)} é chamada por cada worker do Pool.
Ela recebe apenas \code{(i\_b, semente\_rodada)} como argumentos (minimiza
o overhead de serialização).

\textbf{Fluxo interno:}
\begin{enumerate}
  \item Cria um gerador aleatório independente com a semente da rodada.
  \item Sorteia os betas dos intervalos.
  \item Calcula as amplitudes de normalização (uma vez por rodada).
  \item Para cada uma das \code{ITERACOES\_PESOS} combinações de pesos:
    calcula o espectro, avalia $R^2$, e se válido, calcula RMSE.
  \item Retorna o dicionário do melhor resultado desta rodada.
\end{enumerate}

Os contadores \code{\_validos}, \code{\_alta\_precisao} e \code{\_total}
são retornados com prefixo \code{\_} para indicar que são auxiliares
(são agregados pelo processo principal mas não salvos no histórico).

\section{Função busca\_duas\_fases()}

Esta função implementa uma estratégia de busca em dois estágios:

\begin{enumerate}
  \item \textbf{Fase 1 (grossa):} usa \code{FASE1\_ITER\_BETA} $\times$
    \code{FASE1\_ITER\_PESOS} combinações com os intervalos originais completos.

  \item \textbf{Cálculo de novos intervalos:} toma os \code{FASE2\_TOP\_REGIOES}
    melhores resultados da fase 1 e calcula:
    \[
      lo_{\text{novo}} = \max(lo_{\text{orig}},\; \min(\beta) - \text{margem})
    \]
    \[
      hi_{\text{novo}} = \min(hi_{\text{orig}},\; \max(\beta) + \text{margem})
    \]
    onde $\text{margem} = \max(\text{span} \times \text{FASE2\_MARGEM},\;
    \text{range} \times 5\%)$.

  \item \textbf{Fase 2 (fina):} busca densa nos novos intervalos reduzidos.
\end{enumerate}

Por padrão, \code{BUSCA\_DUAS\_FASES = False} — a GUI usa apenas busca
simples com NM.

\section{Geração e Salvamento de Gráficos}

A função \code{gerar\_e\_salvar\_grafico()} cria uma figura de 11$\times$7
polegadas com dois painéis via \code{matplotlib.gridspec}:

\begin{description}[leftmargin=2.5em, style=nextline]
  \item[\textbf{Painel superior (razão 3:1)}]
    Escala log-log com os dados experimentais (pontos pretos) e as curvas
    teóricas: Total (vermelho sólido), MB (roxo), EXP (azul), G1 (verde),
    G2 (laranja). Título com R², RMSE, betas e equação da reta de resíduos.

  \item[\textbf{Painel inferior}]
    Resíduos relativos $E_i = (S_{\text{teo}} - S_{\text{exp}})/S_{\text{exp}}$
    como barras coloridas (vermelho = positivo, azul = negativo) com a reta
    de tendência sobreposta.
\end{description}

O painel lateral com tabela de parâmetros é implementado em \code{\_painel\_parametros()},
que usa \texttt{ax.text()} com coordenadas normalizadas (não \texttt{matplotlib.table}
que tem problemas de layout com fontes monospace).

\section{Modo de Validação Manual}

O modo \code{MODO\_VALIDACAO = True} permite calcular e visualizar o espectro
para um conjunto de parâmetros conhecido, sem rodar a otimização. Útil para:
\begin{itemize}
  \item Verificar se os parâmetros publicados reproduzem o resultado.
  \item Usar como ponto de partida para um refinamento NM dirigido
    (\code{VALIDACAO\_REFINAR\_NM = True}).
  \item Diagnóstico de problemas com a planilha de entrada.
\end{itemize}

Os parâmetros são definidos em \code{VALIDACAO\_PARAMS}, que contém todos
os betas, energias e pesos a serem testados.


% =============================================================================
\chapter{NT\_app\_ctk.py — Interface Gráfica}
\label{cap:gui}
% =============================================================================

\section{Arquitetura Geral}

A GUI é construída com \textbf{customtkinter}, uma biblioteca que adiciona
dark theme nativo e widgets modernos ao Tkinter padrão. A classe principal
\code{NTApp} herda de \code{ctk.CTk} e organiza a interface em:

\begin{itemize}
  \item \textbf{Header fixo}: título da aplicação com logo LMS e subtítulo.
  \item \textbf{Barra de experimento}: campo para caminho do arquivo com botão
    ``Procurar'' e aviso de pasta sincronizada (OneDrive/Dropbox).
  \item \textbf{6 abas principais}: Início, Configuração, Sequência, Executar,
    Resultados, Histórico.
\end{itemize}

A paleta de cores é derivada dos tons do GitHub Dark:
\code{BG = "\#0e1117"} (fundo principal), \code{AZUL = "\#388bfd"} (destaque),
\code{VERDE = "\#238636"} (sucesso), \code{VERM = "\#b91c1c"} (erro).

\section{As Seis Abas}

\subsection{Aba Início}

Apresentação visual com passos numerados (``Como começar'') e cards
descrevendo os dois formatos de entrada aceitos. É a primeira tela que
o usuário vê e serve como guia rápido de uso.

\subsection{Aba Configuração}

Dividida em duas colunas:

\textbf{Coluna esquerda:}
\begin{description}[leftmargin=2.5em, style=nextline]
  \item[Seção ``Busca aleatória'']
    Campos: Iterações beta, Iterações pesos, R² mínimo, RMSE alvo,
    Semente, slider de núcleos CPU.
  \item[Seção ``Nelder-Mead'']
    Checkbox de ativação, Max iter/candidato, Máx reinícios, slider de
    tolerância ($10^{-15}$ a $10^{-3}$), checkbox ``Usar clusters'',
    Limiar cluster, Máx clusters, Top candidatos.
\end{description}

\textbf{Coluna direita:}
\begin{description}[leftmargin=2.5em, style=nextline]
  \item[Seção ``Curvas ativas'']
    Checkboxes para MB, Exponencial, G1, G2, G3, e campo de peso mínimo MB.
  \item[Seção ``Energias das gaussianas'']
    Para cada gaussiana: campo de energia em eV (vazio = usa planilha),
    checkbox ``NM livre'' (permite ao NM otimizar a energia),
    checkbox ``Gravar'' (visível apenas no formato CSV — grava no \texttt{parametros.toml}).
  \item[Seção ``Intervalos de beta'']
    Tabela com Min/Max para cada curva variável (EXP, G1, G2, G3).
  \item[Seção ``Saídas'']
    Campo para número de gráficos PNG a salvar.
\end{description}

Cada seção tem um botão ``?'' que abre um popup de ajuda contextual
(\code{help\_popup()}).

\subsection{Aba Sequência}

Permite enfileirar múltiplas planilhas para processamento automático.
Quando ativa (\code{v\_usar\_seq = True}), a planilha da barra superior
é ignorada e o código processa cada item da fila com as curvas ativas
definidas por item.

O botão ``Adicionar'' cria entradas na lista com o caminho do arquivo e
as curvas selecionadas. Cada entrada tem um botão ``✕'' para remoção.

\subsection{Aba Executar}

Contém:
\begin{itemize}
  \item Botão verde ``Iniciar análise'' e botão vermelho ``Parar''.
  \item Barra de progresso com percentual extraído das linhas de log via
    regex: \code{r'\textbackslash(\textbackslash s*([\textbackslash d.]+)\%\textbackslash)'}.
  \item Terminal monoespaçado com saída em tempo real do subprocesso.
\end{itemize}

\subsection{Aba Resultados}

Exibe automaticamente ao final de uma execução: nome da pasta, R² e RMSE
do melhor resultado (extraídos do arquivo \texttt{.txt}), e botões para
abrir a pasta e o Excel.

\subsection{Aba Histórico}

Lista todas as execuções anteriores encontradas em \texttt{Resultados/}.
Cada entrada é expansível (botão ``runs'') para mostrar as rodadas individuais
válidas com seus parâmetros. Checkboxes em dois níveis:

\begin{itemize}
  \item \textbf{Checkbox principal}: seleciona o melhor resultado global do experimento.
  \item \textbf{Checkboxes de runs}: seleciona rodadas individuais (NM1, NM2, etc.).
\end{itemize}

Com 2 ou mais selecionados, o botão ``Comparar'' fica ativo.

\section{Geração do Script Temporário}

Este é o mecanismo central de execução. A função \code{criar\_script\_temporario()}:

\begin{enumerate}
  \item Lê \arquivo{analise\_natalia\_time.py} inteiro do disco.
  \item Localiza os marcadores \texttt{SEÇÃO 2} e \texttt{SEÇÃO 3} com regex:
    \begin{lstlisting}[caption={Localização dos marcadores de seção}]
m2s = re.search(r'#\s*={10,}\s*\n#\s*SECAO 2', codigo)
m3s = re.search(r'#\s*={10,}\s*\n#\s*SECAO 3', codigo)
    \end{lstlisting}
  \item Substitui o bloco entre SEÇÃO 2 e SEÇÃO 3 pelo texto gerado por
    \code{gerar\_config\_python(cfg)}, que serializa todos os parâmetros
    da GUI como código Python válido.
  \item Salva em um arquivo temporário em \texttt{tempfile.mkstemp()}.
\end{enumerate}

O design de injeção em vez de passagem por argumentos foi escolhido porque
o motor de análise foi originalmente projetado para execução independente
com configuração estática. Isso permite que qualquer edição em
\arquivo{analise\_natalia\_time.py} tome efeito imediatamente, sem precisar
rebuildar o executável.

\section{Execução como Subprocess}

A execução ocorre em uma thread separada para não bloquear a GUI:

\begin{lstlisting}[caption={Execução como subprocess em thread daemon}]
def _runner():
    cmd = [sys.executable, "-u", tmp]
    proc = subprocess.Popen(
        cmd,
        stdout=subprocess.PIPE, stderr=subprocess.STDOUT,
        text=True, encoding="utf-8", errors="replace",
        bufsize=1, env=env,
    )
    self.proc_ref = proc
    for line in proc.stdout:
        self._log_q.put(line.rstrip())
        # Captura caminho da pasta de resultados da saída
        for token in line.split():
            if os.path.isdir(token) and "Resultados" in token:
                ultima_pasta = token

threading.Thread(target=_runner, daemon=True).start()
self.after(300, self._poll)
\end{lstlisting}

O método \code{\_poll()} é chamado a cada 300~ms e processa a fila
\code{\_log\_q} para atualizar o terminal e a barra de progresso,
sem bloquear o event loop do Tkinter.

\section{Mecanismo de Pause/Stop}

O botão ``Parar'' da GUI:
\begin{enumerate}
  \item Desabilita o próprio botão para evitar duplo clique.
  \item Identifica a pasta de resultados atual (lida do output do processo).
  \item Cria o arquivo \code{PARAR.txt} nessa pasta.
  \item Inicia uma thread auxiliar que aguarda até 10 segundos e, se
    o processo não encerrar, chama \code{proc.terminate()}.
\end{enumerate}

\section{Janela de Comparação}

Acessível pela aba Histórico, a janela \code{\_janela\_comparacao()} exibe:

\begin{enumerate}
  \item \textbf{Tabela de parâmetros}: grade com R², RMSE, pesos e betas de
    cada seleção lado a lado, com linhas alternadas cinza/escuro.

  \item \textbf{Gráfico de barras empilhadas horizontais}: mostra a
    contribuição percentual de cada componente (MB, EXP, G1, G2, G3) para
    cada experimento/rodada. Usa matplotlib embutido via
    \code{FigureCanvasTkAgg}.

  \item \textbf{Grade de PNGs}: exibe os gráficos de cada rodada em grade
    de 2 colunas. A função \code{\_encontrar\_png()} localiza os arquivos:
    \begin{lstlisting}[caption={Localização de PNGs por rodada}]
rod_fmt = str(int(float(rod_str))).zfill(4)
pngs = glob.glob(
    os.path.join(pasta, "graficos_rodadas",
                 f"rod{rod_fmt}_*.png"))
    \end{lstlisting}
    Para rodadas NM, usa o nome diretamente (``NM1'', ``NM2'', etc.)
    sem formatação numérica.
\end{enumerate}


% =============================================================================
\chapter{Migrador NT/migrador.py — Conversor de Formatos}
\label{cap:migrador}
% =============================================================================

\section{Propósito}

O Migrador é um utilitário standalone que converte dados experimentais do
formato original (planilha Excel com estrutura específica) para o formato
CSV+TOML consumido pelo Natalia Time. Permite que o usuário trabalhe sem
ter o Excel instalado na máquina de análise.

\section{Fluxo de Uso}

\begin{enumerate}
  \item Usuário informa a \textbf{pasta raiz de destino} e o \textbf{nome do experimento}.
  \item Preenche os \textbf{parâmetros físicos}: temperatura do gás, massas,
    parâmetros geométricos do espectrômetro e offset de tempo.
  \item Fornece os dados via:
    \begin{description}
      \item[\textbf{Aba Automático}] Seleciona o arquivo \texttt{.xlsx} de dados brutos;
        o Migrador lê a aba \texttt{Analise\_tempo} nas colunas L e M.
      \item[\textbf{Aba Manual}] Cola os dados copiados diretamente do Excel
        (colunas separadas por tab ou vírgula).
    \end{description}
  \item Clica em ``Criar experimento''.
\end{enumerate}

O resultado é uma pasta \texttt{<destino>/<nome>/} contendo \texttt{dados.csv}
e \texttt{parametros.toml}.

\section{Parâmetros de Entrada}

\begin{table}[h!]
  \centering
  \caption{Parâmetros configuráveis no Migrador}
  \begin{tabular}{llp{6cm}}
    \toprule
    \textbf{Campo} & \textbf{Padrão} & \textbf{Descrição} \\
    \midrule
    \code{T\_gas} & 300 & Temperatura do gás de expansão (K) — usada para calcular $\beta_{\text{mb}}$ \\
    \code{m\_frag} & 16,0 & Massa do fragmento iônico (ex: 16 para O$^+$) \\
    \code{m\_mol} & 32,0 & Massa da molécula-mãe (ex: 32 para O$_2$) \\
    \code{DE} & 350,0 & Campo de extração em eV \\
    \code{D} & 6,8 & Diâmetro do colimador em mm \\
    \code{LL} & 6,8 & Comprimento efetivo do cilindro de ionização em mm \\
    \code{I7\_NORM} & 500,0 & Tempo de referência para normalização (ns) \\
    \code{t\_offset} & 600,0 & Offset de tempo eletrônico a somar nos tempos brutos (ns) \\
    \bottomrule
  \end{tabular}
\end{table}

\subsection{Cálculo Automático de beta\_mb}

O parâmetro $\beta_{\text{mb}}$ não é inserido manualmente — é calculado
pela fórmula:
\begin{equation}
  \beta_{\text{mb}} = \sqrt{\frac{4{,}04 \times 10^5 \cdot \frac{300}{T_{\text{gas}}} \cdot m_{\text{mol}}}{2}}
\end{equation}
e gravado diretamente no \texttt{parametros.toml}. Isso elimina a necessidade
do usuário calcular e inserir $\beta_{\text{mb}}$ manualmente.

\section{Parsing de Dados}

\subsection{Modo Automático — Leitura do .xlsx}

A função \code{\_ler\_xlsx()} lê especificamente:
\begin{itemize}
  \item Aba: \texttt{Analise\_tempo}
  \item Coluna L (índice 12): tempo em ns
  \item Coluna M (índice 13): sinal normalizado
  \item A partir da linha 13 (\code{LINHA\_DADOS = 13})
\end{itemize}

Pontos com $t = 0$ e $s \approx 1{,}0$ são filtrados (são pontos de referência).

\subsection{Modo Manual — Cola do Excel}

A função \code{\_parsear\_dados\_manual()} aceita:
\begin{itemize}
  \item \textbf{Tab-separado}: resultado de copiar células do Excel (padrão).
  \item \textbf{Vírgula/ponto-e-vírgula}: separadores alternativos.
  \item Vírgulas decimais são convertidas para pontos antes do \code{float()}.
\end{itemize}

O widget de texto exibe um placeholder com exemplo de formato enquanto vazio.
A contagem de pontos válidos é atualizada em tempo real a cada tecla pressionada.

\section{Aplicação do Offset de Tempo}

O offset de tempo corrige o atraso eletrônico da aquisição: os dados brutos
têm tempos menores do que os tempos físicos reais de voo. Ao criar o CSV,
o Migrador soma \code{t\_offset} a cada tempo:

\begin{lstlisting}[caption={Aplicação do offset ao escrever dados.csv}]
with open(os.path.join(pasta_exp, "dados.csv"), "w") as f:
    f.write("tempo,sinal\n")
    for t, s in pontos:
        f.write(f"{t + t_offset},{s}\n")
\end{lstlisting}

O valor padrão de 600~ns é um ponto de partida — deve ser ajustado pelo
usuário conforme o equipamento. O offset é registrado como comentário
no cabeçalho do \texttt{parametros.toml} para rastreabilidade.

\section{Estrutura dos Arquivos Gerados}

\subsection{dados.csv}

\begin{verbatim}
tempo,sinal
610.0,1.0
660.0,0.985348
710.0,0.950200
...
\end{verbatim}
(Os tempos brutos de 10, 60, 110, ... ns são somados com o offset de 600 ns.)

\subsection{parametros.toml}

\begin{lstlisting}[language={},caption={Estrutura do parametros.toml gerado}]
# Parametros fisicos do experimento
# Gerado pelo Natalia Time - Criador de Experimento
# T_gas = 300.0 K  ->  beta_mb = 2572.630956
# Offset de tempo aplicado: 600.0 ns

# Parametros da molecula
beta_mb = 2572.630956
m_frag  = 16.0
m_mol   = 32.0

# Parametros geometricos do espectrometro DETOF
DE      = 350.0
D       = 6.8
LL      = 6.8
I7_NORM = 500.0

# Energias das gaussianas (eV)
# Definidas e gravadas pelo Natalia Time
# e_g3 = 0.7
# e_g4 = 3.0
# e_g5 = 2.0
\end{lstlisting}

As energias ficam comentadas porque dependem do experimento específico
e devem ser determinadas pelo físico antes da análise.

\section{Validações e Tratamento de Erros}

O Migrador valida:
\begin{itemize}
  \item Campos obrigatórios não podem estar vazios.
  \item Todos os parâmetros numéricos devem ser válidos (aceita vírgula decimal).
  \item Nome do experimento não pode conter caracteres inválidos para nome de pasta
    (\texttt{\textbackslash~/:*?"<>|}).
  \item Mínimo de 5 pontos experimentais válidos.
  \item Se a pasta de destino já existir, solicita confirmação antes de sobrescrever.
\end{itemize}


% =============================================================================
\chapter{Portabilidade e Distribuição}
\label{cap:distribuicao}
% =============================================================================

\section{Dependências Python}

\begin{table}[h!]
  \centering
  \caption{Dependências Python com versões recomendadas}
  \begin{tabular}{llp{6cm}}
    \toprule
    \textbf{Pacote} & \textbf{Versão recomendada} & \textbf{Uso} \\
    \midrule
    \code{numpy}         & $\geq 1.24$ & Álgebra linear, grade de velocidades, einsum \\
    \code{scipy}         & $\geq 1.10$ & Nelder-Mead, clustering hierárquico \\
    \code{matplotlib}    & $\geq 3.7$  & Gráficos, viewer interativo \\
    \code{pandas}        & $\geq 2.0$  & DataFrames para resultados \\
    \code{openpyxl}      & $\geq 3.1$  & Leitura e escrita de .xlsx/.xlsm \\
    \code{customtkinter} & $\geq 5.2$  & GUI com dark theme \\
    \code{darkdetect}    & $\geq 0.8$  & Detecção de tema do sistema (dep. do customtkinter) \\
    \code{tomllib}       & embutido (Python $\geq 3.11$) & Leitura de parametros.toml \\
    \code{tkinter}       & embutido    & Base para GUI e para o Migrador \\
    \bottomrule
  \end{tabular}
\end{table}

\textbf{Nota:} \code{tomllib} é parte da biblioteca padrão do Python 3.11+.
Para Python 3.10 ou anterior, seria necessário usar \code{tomli} como
dependência externa.

\section{Build com PyInstaller no Windows}

\subsection{NataliaTime.exe}

O arquivo \arquivo{NataliaTime.spec} define o build do executável principal.
O ponto de entrada é \arquivo{NT\_app\_ctk.py}. A instrução essencial é:

\begin{lstlisting}[language={},caption={Trecho relevante do NataliaTime.spec}]
from PyInstaller.utils.hooks import collect_data_files
datas = collect_data_files("customtkinter")
\end{lstlisting}

Isso inclui os arquivos de tema do customtkinter (JSON, imagens) que são
necessários para o dark mode funcionar no executável.

\textbf{Importante}: \arquivo{analise\_natalia\_time.py} \textbf{não} é
embutido no executável — ele é lido do disco em tempo de execução via o
mecanismo de script temporário. Isso significa que:
\begin{itemize}
  \item Mudanças em \arquivo{analise\_natalia\_time.py} entram em vigor
    imediatamente, sem rebuild.
  \item O arquivo \arquivo{analise\_natalia\_time.py} \textbf{deve} estar na
    mesma pasta que \texttt{NataliaTime.exe}.
\end{itemize}

O script \arquivo{build\_nataliatime.py} automatiza o processo:
\begin{lstlisting}[caption={Build automatizado}]
# python build_nataliatime.py
# Limpa build/ e dist/, depois roda PyInstaller com NataliaTime.spec
\end{lstlisting}

\subsection{Migrador.exe}

Construído de forma análoga a partir de \arquivo{Migrador NT/migrador.py}
usando \arquivo{Migrador NT/Migrador.spec} e \arquivo{build\_migrador.py}.
O Migrador não tem dependências externas pesadas, resultando em um executável
menor.

\section{Pacote de Distribuição}

O pacote de distribuição é \texttt{pacote\_distribuicao/NataliaTime\_v1.1.0.zip}
($\approx$122 MB) e contém:

\begin{description}[leftmargin=2.5em, style=nextline]
  \item[\texttt{NataliaTime.exe}] Executável principal (requer as dependências ao lado).
  \item[\texttt{Migrador.exe}] Utilitário de conversão de formatos.
  \item[\texttt{NT\_app\_ctk.py}] Interface gráfica (não embutida no exe).
  \item[\texttt{analise\_natalia\_time.py}] Motor de análise (lido do disco em runtime).
  \item[\texttt{LEIA-ME.txt}] Guia do usuário em português.
\end{description}

\subsection{Instalação para o Usuário Final}

\begin{enumerate}
  \item Descompactar o ZIP em qualquer pasta (preferencialmente fora do OneDrive).
  \item Executar \texttt{NataliaTime.exe} diretamente — sem instalação.
  \item Para converter planilhas antigas, usar \texttt{Migrador.exe}.
\end{enumerate}

\section{Portabilidade para Linux}

A portabilidade para Linux é possível mas requer as seguintes adaptações:

\begin{description}[leftmargin=2.5em, style=nextline]
  \item[\textbf{GUI}]
    customtkinter funciona no Linux via Tkinter, mas pode exigir instalação
    de pacotes do sistema (\code{python3-tk}, \code{tcl-tk-dev}) e ajustes
    de fonte para Segoe UI (não nativa no Linux — substituir por DejaVu Sans
    ou Liberation Sans).

  \item[\textbf{Abertura de arquivos}]
    O código usa \code{os.startfile(path)} para abrir pastas e Excel, que é
    exclusivo do Windows. No Linux, substituir por:
    \begin{lstlisting}[caption={Substituição de os.startfile() para Linux}]
import subprocess
subprocess.Popen(["xdg-open", path])
    \end{lstlisting}

  \item[\textbf{Separador de caminhos}]
    O código já usa \code{os.path.join()} na maior parte, mas há
    strings com caminhos hardcoded usando \texttt{\textbackslash} que
    precisariam ser corrigidas.

  \item[\textbf{Ícone e título da janela}]
    A chamada \code{ctypes.windll.shell32.SetCurrentProcessExplicitAppUserModelID()}
    no Migrador precisa ser envolvida em \code{try/except} (já está).

  \item[\textbf{PyInstaller no Linux}]
    Produz binário ELF nativo. O processo de build é idêntico ao Windows
    (\code{python build\_nataliatime.py}) mas requer a dependência
    \code{binutils} no sistema.

  \item[\textbf{tomllib}]
    Disponível nativamente no Python 3.11+ em todas as plataformas. Sem
    mudanças necessárias.
\end{description}

O motor de análise (\arquivo{analise\_natalia\_time.py}) é puramente
numérico e não usa nenhuma API específica do Windows — ele rodaria no Linux
sem modificações.


% =============================================================================
\appendix
\chapter{Tabela de Parâmetros de Configuração (Seção 2)}
\label{ap:config}
% =============================================================================

A tabela a seguir lista todos os parâmetros configuráveis da Seção~2 de
\arquivo{analise\_natalia\_time.py}, seus valores padrão e a descrição.
Estes são os parâmetros injetados pela GUI ao gerar o script temporário.

\begin{longtable}{lll}
  \caption{Parâmetros da Seção 2 — Configuração} \\
  \toprule
  \textbf{Parâmetro} & \textbf{Padrão} & \textbf{Descrição} \\
  \midrule
  \endfirsthead
  \toprule
  \textbf{Parâmetro} & \textbf{Padrão} & \textbf{Descrição} \\
  \midrule
  \endhead
  \bottomrule
  \endfoot
  \code{SCORE\_MINIMO}       & 0,99  & R² mínimo para resultado válido \\
  \code{DESVIO\_ALVO}        & 0,04  & RMSE máximo para ``alta precisão'' \\
  \code{ITERACOES\_BETA}     & 10    & Combinações de beta sorteadas \\
  \code{ITERACOES\_PESOS}    & 10    & Pesos testados por beta \\
  \code{SEMENTE}             & None  & Semente aleatória (None = não reproduzível) \\
  \code{N\_PROCESSOS}        & None  & Núcleos CPU (None = todos) \\
  \code{ARQUIVO\_PARAR}      & PARAR.txt & Nome do arquivo de sinal de parada \\
  \code{INTERVALOS\_BETA}    & vide código & Intervalos $[lo, hi]$ para cada beta \\
  \code{CURVAS\_PESO\_ATIVAS} & vide código & Quais curvas participam da otimização \\
  \code{ENERGIAS\_VALOR}     & todos None & Valores fixos para energias (None = planilha) \\
  \code{ENERGIA\_NM\_LIVRE}  & todos False & Permite NM ajustar a energia \\
  \code{INTERVALOS\_ENERGIA} & (0,1; 6,0) & Limites para energias livres no NM \\
  \code{PESO\_MIN}           & p\_mb=0,01 & Peso mínimo por curva \\
  \code{SALVAR\_GRAFICOS\_TOP} & 20  & Máximo de PNGs a salvar \\
  \code{VIEWER\_TOP}         & 50   & Máximo no viewer interativo \\
  \code{USAR\_NELDER\_MEAD}  & True & Ativa refinamento local \\
  \code{NM\_TOP\_CANDIDATOS} & 5    & Candidatos ao NM (sem clusters) \\
  \code{NM\_MAX\_ITER}       & 2000 & Iterações máximas por candidato \\
  \code{NM\_TOLERANCIA}      & $10^{-9}$ & Critério de convergência \\
  \code{NM\_MAX\_REINICIAR}  & 3    & Reinícios por candidato \\
  \code{NM\_PERTURB\_ESCALA} & 0,10 & Intensidade da perturbação no reinício \\
  \code{NM\_USAR\_CLUSTERS}  & True & Seleção por clustering hierárquico \\
  \code{NM\_CLUSTER\_LIMIAR} & 0,25 & Distância de corte do clustering \\
  \code{NM\_CLUSTERS\_MAX}   & 10   & Máximo de clusters (candidatos) \\
  \code{BUSCA\_DUAS\_FASES}  & False & Ativa busca em dois estágios \\
  \code{ADAPTAR\_INTERVALOS} & False & Ativa adaptação automática de intervalos \\
  \code{MODO\_VALIDACAO}     & False & Ativa modo de validação manual \\
  \code{VALIDACAO\_REFINAR\_NM} & False & Refina com NM após validação manual \\
\end{longtable}


\chapter{Glossário de Termos Técnicos}
\label{ap:glossario}

\begin{description}[leftmargin=2.5em, style=nextline]
  \item[\textbf{Beta ($\beta$)}]
    Parâmetro de largura das distribuições de energia cinética. Valores
    maiores correspondem a distribuições mais estreitas (fragmentos mais
    ``frios''). Dimensional: inverso de velocidade.

  \item[\textbf{DETOF}]
    \textit{Drift-time Electrostatic Time-of-Flight}: espectrômetro de
    massa baseado no tempo de deriva de íons após aceleração eletrostática.

  \item[\textbf{Dirichlet}]
    Distribuição de probabilidade sobre o simplex $K$-dimensional. Usada
    para amostrar vetores de pesos que somam 1 de forma uniforme.

  \item[\textbf{Gauss-Hermite}]
    Método de quadratura numérica para integrais da forma
    $\int_{-\infty}^{+\infty} f(x) e^{-x^2} dx$. O código usa 4 pontos
    (suficiente para as integrais envolvidas).

  \item[\textbf{Logit/softmax}]
    Transformação que mapeia $K$ valores reais irrestringidos em $K$ probabilidades
    positivas que somam 1. Usada para parametrizar os pesos no Nelder-Mead.

  \item[\textbf{Nelder-Mead}]
    Método de otimização local sem gradiente baseado em um simplex geométrico
    que se deforma iterativamente. Robusto para funções com ruído numérico.

  \item[\textbf{R² (coeficiente de determinação)}]
    Medida de qualidade do ajuste entre 0 e 1. $R^2 = 1$ indica ajuste perfeito.
    No contexto deste código, $R^2 \geq 0{,}99$ é o critério de validade.

  \item[\textbf{RMSE}]
    \textit{Root Mean Square Error}: raiz do erro quadrático médio.
    Neste código, calcula o RMSE dos resíduos relativos após subtração de
    tendência linear.

  \item[\textbf{Softmax}]
    Ver Logit/softmax.

  \item[\textbf{Simplex (Nelder-Mead)}]
    Poliedro $N$-dimensional com $N+1$ vértices usado pelo método de
    Nelder-Mead para explorar o espaço de parâmetros.

  \item[\textbf{TOML}]
    \textit{Tom's Obvious Minimal Language}: formato de arquivo de
    configuração legível por humanos, usado no \texttt{parametros.toml}.

  \item[\textbf{Worker (multiprocessing)}]
    Processo filho criado pelo Pool. Cada worker recebe um subconjunto
    das rodadas e as processa independentemente em paralelo.
\end{description}

\end{document}
